\section{The cylinder exponent}\label{sec:E-cylinder-exponent}

We can now sum the contour sectors. The interpolation just proved is
needed only during the summation over imbalance: it gives a holomorphic
function approximating the fixed-dimension integrals to arbitrarily high
exponential precision. After moving that sum, we restore the scalar
normalization and obtain the physical lower bound from positivity.


Take the balanced physical cylinder with \(N=2m\). Its row
parameters are \(2\lambda\) on the first group and \(\lambda\)
on the second. Their multiplicative site values \(q^2,q\) become
\(1,q^{-1}\) under the common rescaling by \(q^{-2}\), which leaves
the contraction unchanged. Thus the logarithmic sites used below are
\(0,-i\beta\), as in Section~\ref{sec:02-cylinder}.

Define the scalar
\begin{equation}\label{eq:04-physical-normalization}
 \Delta_N=
 \frac{\widetilde h_N^2}
 { (\gamma/2)^N\prod_{i<j}H_*(a_i-a_j)^2},
 \qquad G_N(\eta)=\Delta_N Z_N(2\cos\eta).
\end{equation}
Theorem~\ref{thm:02-denominator} proves that \(\Delta_N\ne0\) and
\(\log|\Delta_N|=(5/24)\log N+o(\log N)\). Dividing the exact
contour-sector identity by the same extracted site product expresses
\(G_N\) as the sum and integral over explicit configurations \(\zeta\)
of a frozen prefactor times
\begin{equation}
 \sum_M e^{i(\eta+\pi f)M}F_{m'(M)}(M).
 \label{eq:G-22}
\end{equation}
Only valid integer indices and dimensions occur. The explicit measures
include their species factorials; the prefactor is independent of
\(M\) and has modulus at most \(\exp(C_hn_e)\).

\begin{proposition}[Summed contour bound]\label{prop:04-sector-sum}
Locally uniformly in the strip \(|\operatorname{Re}\eta|<\pi\),
\begin{equation}
 \log |G_N(\eta)|\le
 3\gamma L/8-2L\,\operatorname{Re}(\eta^2)/(4\kappa)+o(L)
 \label{eq:G-23}
\end{equation}
\end{proposition}

\begin{proof}
Fix \(|\operatorname{Re}\eta|\le a_0<\pi,\ |\operatorname{Im}\eta|\le b_0\). Choose the small parameters for \eqref{eq:G-16},\eqref{eq:G-17} as in \eqref{eq:G-6},\eqref{eq:G-7}.
We first truncate the damped real sectors and replace their integrals by
the holomorphic approximation. A rectangle summation formula then
separates the integer sum into Fourier aliases, each of which can be
moved to the height selected by its trial bound.

\paragraph{Truncation and holomorphic replacement.}

First the untruncated real bound in part 1 of
Proposition~\ref{prop:04-integer-bounds} lets us restrict \(n_e,\ell_e\le L^3\) and \(|M|<T\), with \(T\) a half-integer of a sufficiently large fixed polynomial scale to be increased below. Indeed \(m'(M)\) on valid indices is \(O(m)\); summing the reserved part of damping in \(M\) costs only \(O(\sqrt L)\). A reserved exponential damping in \(\ell_e\) makes the summed explicit factorial measures, even with the per-particle factors, cost \(\exp(O_h(L))\), just by integration of each variable. Thus these tail errors after prefactors are, say, \(O(\exp(-c_h L^2))\). Under the resulting restrictions all integers inside the \(M\)-cut are valid for large \(L\); \(\nu=m'(M)\) lies within a polynomial distance of \(m\).

At each retained \(\zeta\) use \eqref{eq:G-20} to replace the integral in \eqref{eq:G-22} by \(H(M)=\widehat F(m'(M),M)\). The precision power and domain power in \eqref{eq:G-20} can be taken sufficiently large after all polynomial bounds below; hence the summed error is negligible, even with the phase, prefactors and integration. For instance even the undamped truncated factorial measures with fixed per-particle factors cost at most \(\exp({\rm poly}(L))\). For each bound \eqref{eq:G-15}–\eqref{eq:G-17} at a particular \(M\) below we can transfer it from all independent integer \(\nu\) near \(m\) to \(H(M)\) by \eqref{eq:G-21}, losing only \(o(L)\) in the log. Here truncate \(j\) polynomially when transferring \eqref{eq:G-16},\eqref{eq:G-17}; in \eqref{eq:G-15} just transfer for \(j=0\) and then insert any alias phase exactly. Indeed the quantitative upper majorants used for \(|F_\nu(M)|\) (after removing phases) are bounded below by \(\exp(-{\rm poly}(L))\) with powers fixed independently of the precision choice. All such transfers are uniform.

\paragraph{The rectangle formula and alias shifts.}
Write \(f_\eta(M)=e^{i(\eta+\pi f)M}H(M)\), and let
\(\mathcal R_T\) be the counterclockwise rectangle with real
endpoints \(\pm T\) and imaginary heights \(\pm1\). Since the
residue of \((1-e^{2\pi iM})^{-1}\) at an integer is
\(-1/(2\pi i)\), the exact summation formula is
\[
 \sum_{k\in\Z,\ |k|<T}f_\eta(k)
       =-\int_{\partial\mathcal R_T}
                \frac{f_\eta(M)}{1-e^{2\pi iM}}\,dM.
\]
At the half-integer vertical sides, the kernel has modulus at most
one. On the two horizontal sides its absolutely convergent expansions
are
\[
 \frac1{1-e^{2\pi iM}}=
 \begin{cases}
 \displaystyle\sum_{j\ge0}e^{2\pi ijM},&\operatorname{Im}M=1,\\[3pt]
 \displaystyle-\sum_{j<0}e^{2\pi ijM},&\operatorname{Im}M=-1.
 \end{cases}
\]
Equation~\eqref{eq:G-15} bounds the alias tails by a geometric
factor \(e^{-2\pi J}\), times an exponential of a fixed polynomial
from the explicit truncations. Thus it suffices to retain
\(|j|\le J=L^{C_J}+O(1)\), choosing \(C_J\) sufficiently large.
This choice can precede the final power of \(T\): the Gaussian
damping in \(\operatorname{Re}M\) is integrable on the whole line,
and even any later fixed polynomial interval length is absorbed by
\(e^{-2\pi J}\).

For each remaining alias, move its horizontal integration at fixed endpoints' real parts to the height \(v\) selected in \eqref{eq:G-6} if \(U>U_0\), or \eqref{eq:G-7} if \(U\le U_0\). The choice depends only on alias, bias, and explicit data, and has a polynomial bound independent of \(T\). The connecting verticals and the original end integrals are all negligible by \eqref{eq:G-15}: there the loss \(c_h T^2/L\) can dominate all other powers, including possible growth from \(|j v|\), by taking the power of \(T\) sufficiently large after \(C_J\). These choices only require a fixed polynomial box in \eqref{eq:G-20}; the precision there is now chosen as just described. All movements of \(H\) times its phases are thus within its holomorphy domain.

\paragraph{Summing the shifted bounds.}
On the shifted lines \eqref{eq:G-16},\eqref{eq:G-17} apply, up to \(o(L)\) by the transfer. In the large-\(U\) case integrate the bound \eqref{eq:G-16} using exponential \(\ell_e\)-damping and factorials, and absorb its \(O_h(L)\) costs by taking \(U_0\) sufficiently large (after \(\delta\) and \(h\)). In the bounded-\(U\) case the bound \eqref{eq:G-17} costs after prefactors and integration at most an additional \(C_{h,\varepsilon}\delta L+o(L)\) in the log, as following \eqref{eq:G-7}. More explicitly this case is supported on explicit configurations with at most \(C_{U_0}\) far particles on \(|u|\le(1-2\delta)L\), no restriction needed for near outer particles other than being exterior. Thus this same integrated bound works for each alias and there are only polynomially many; likewise horizontal length costs at most polynomially. This pointwise split does not require any contour movement of explicit variables. All constants allow \(\delta\) small after \(\varepsilon,h\), then \(U_0\) arbitrarily large fixed, little-oh errors then for these fixed choices. Sending the small errors to zero proves \eqref{eq:G-23}.
\end{proof}

\begin{proof}[Proof of Theorem~\ref{E:cylinder-exponent}]
By \eqref{eq:04-physical-normalization}, the physical exponent is the
contour numerator exponent minus the scalar exponent. The deformation
determinant contributes \(7/8\), the two reference masses contribute
\(-1/2\), and the scalar normalization contributes \(-5/24\), all
as coefficients of \(\log N\). Thus the leading constants are
\[
 \frac78-\frac12-\frac5{24}=\frac16,
 \qquad \frac1{2\gamma\kappa}=\frac2{3\pi^2},
\]
since \(\gamma=4/3\) and \(\kappa=(3\pi/4)^2\). Thus
Proposition~\ref{prop:04-sector-sum} gives, locally uniformly in
\(|\operatorname{Re}\eta|<\pi\),
\[
 u_N(\eta):=\frac{\log|Z_N(2\cos\eta)|}{\log N}
 \le\mathcal H(\eta)+o(1),\qquad
 \mathcal H(\eta)=\frac16-\frac{2}{3\pi^2}
                                  \operatorname{Re}(\eta^2).
\]
\paragraph{Propagation from the empty-configuration value.}
The function \(\mathcal H\) is harmonic. Since \(Z_N\) has
nonnegative coefficients and \(Z_N(0)=1\), the interior point
\(\eta=\pi/2\) satisfies \(u_N(\pi/2)=\mathcal H(\pi/2)=0\).
We show that every fixed disk inside the strip contains points
\(\eta_N\) with
\(u_N(\eta_N)\ge\mathcal H(\eta_N)-o(1)\).

Join the anchor to the disk by a compact path in the strip and
choose a connected neighborhood with positive clearance from its
boundary. The locally uniform upper bound permits a constant
\(\epsilon_N\downarrow0\) such that
\(v_N=\mathcal H+\epsilon_N-u_N\ge0\) on this neighborhood.
It is superharmonic and locally integrable; zeros of \(Z_N(2\cos\eta)\)
give only logarithmic singularities, and the anchor ensures this
holomorphic function is not identically zero. At the anchor,
\(v_N=\epsilon_N\).

Choose a finite chain of disks of radius \(r\), with successive
centers at distance at most \(r\), leading into the target disk.
Take \(r\) small enough that, from any point in one disk, the
radius-\(3r\) disk stays in the chosen neighborhood and covers the
next disk. The area mean inequality for a nonnegative superharmonic
function implies that if \(v_N(z)\le a_N\) at the chosen point,
then its integral over the next radius-\(r\) disk is at most
\(9\pi r^2a_N\). Some point there therefore has value at most
\(9a_N\). Starting at the anchor and iterating the fixed finite
chain yields \(v_N(\eta_N)\le9^k\epsilon_N=o(1)\) in the target.

\paragraph{Comparison with positive real fugacities.}
The disk argument supplies a near-saturating point, not a pointwise
lower bound throughout the disk. Positivity of the physical polynomial
converts this into the desired lower bound at a positive fugacity.
For \(\chi=2\), take target disks about small positive real values
of \(\eta\), sufficiently narrow that \(|2\cos\eta|<2\)
throughout. Positivity of coefficients gives
\(Z_N(2)\ge|Z_N(2\cos\eta_N)|\). Thus the lower limit of
\(\log Z_N(2)/\log N\) is at least the infimum of
\(\mathcal H\) on each such disk. Letting these disks approach
zero gives \(1/6\); the upper bound at zero matches it.

If \(0<\chi<2\), put \(\eta_\chi=\arccos(\chi/2)\).
Choose real disk centers decreasing to \(\eta_\chi\) from above,
and radii sufficiently small that \(|2\cos\eta|\le\chi\)
throughout. The same positivity argument gives the lower exponent
\(\mathcal H(\eta_\chi)\), matching the upper bound there.
If \(\chi>2\), use
\(\eta_\chi=i\operatorname{arcosh}(\chi/2)\) and disk centers
approaching it along the imaginary axis from smaller imaginary
parts. Again their radii can be chosen so that
\(|2\cos\eta|\le\chi\). This proves the remaining formula.
The argument obtains a lower bound from the polynomial's positivity
and the interior anchor; it requires no lower estimate on a single
complex contour sector.
\end{proof}
